\section{Evaluation}
\label{sec:eval}

We evaluate three claims. First, that the structural prior (the per-category $L_C$) carries a measurable amount of useful knowledge across same-category apps. Second, that online evolution of LAYER~2 on the target app's small adaptation set adds further useful knowledge on top. Third, that joint weight and context adaptation completes the picture by lifting performance on apps where the policy itself, not just its prompt, needs to specialize. The first two claims are tested by the B1$\to$B2 and B2$\to$B3 deltas, the third by B3$\to$B4. Throughout, the Tier-C far-transfer panel acts as a built-in null control on the injection mechanism.

\subsection{Benchmark and splits}
\label{sec:eval-benchmark}

The benchmark is built on AndroidWorld~\citep{rawles2024androidworld}, extended with six apps from BMOCA (Calculator, Snapseed, Wikipedia) and AndroidLab~\citep{xu2024androidlab} (Bluecoins, Maps.me, Pi Music) for 25 apps in total. The apps cover 12 Google Play Store categories and 192 app-attributable task templates. We split them into three pools as in Table~\ref{tab:splits}. The source pool spans the six Play Store categories that have at least two apps; the Tier-B target pool is composed of apps in categories that are represented in the source pool (near-transfer); the Tier-C target pool draws from single-app categories absent from the source pool (far-transfer).

\begin{table}[t]
\small
\centering
\begin{tabular}{lrrr}
\hline
\textbf{Pool} & \textbf{Apps} & \textbf{Templates} & \textbf{Episodes} \\
\hline
Source pool                & 12 & 96  & 480 (K=5) \\
Tier-B $T_{\text{adapt}}$  &  6 & 32  & 160 (K=5) \\
Tier-B $T_{\text{eval}}$   &  6 & 18  &  54 (K=3) \\
Tier-C $T_{\text{adapt}}$  &  7 & 29  & 145 (K=5) \\
Tier-C $T_{\text{eval}}$   &  6 & 17  &  51 (K=3) \\
\hline
\textbf{Total}             & 25 & 192 & 890 \\
\hline
\end{tabular}
\caption{Pool composition. Templates are schema-level units; episodes count templates times K seeds. Tier-C contains one extra app (\texttt{simple\_draw\_pro}) with $N{=}1$ template assigned entirely to $T_{\text{adapt}}$; its $T_{\text{eval}}$ is empty and the app is excluded from the Tier-C aggregate.}
\label{tab:splits}
\end{table}

Within every target app, templates are partitioned 60/40 into $T_{\text{adapt}}$ and $T_{\text{eval}}$ by deterministic alphabetical ordering. The two splits are \emph{template-disjoint}: the templates themselves differ, so $T_{\text{eval}}$ measures within-app generalization of the adaptation procedure rather than parameter robustness on memorized templates---the support/query convention of meta-learning~\citep{finn2017maml} and the train-level / test-level convention of Procgen~\citep{cobbe2020procgen}. Seeds also differ: $T_{\text{adapt}}$ uses $\{30, 31, 32, 33, 34\}$ ($K{=}5$), and $T_{\text{eval}}$ uses $\{40, 41, 42\}$ ($K{=}3$). Stacked together, the three independent splits---pool, category, and template---give an honest measurement of cross-app and cross-category generalization. Absolute numbers under this protocol are lower than the same models' AndroidWorld leaderboard numbers by design.

\subsection{Baselines}
\label{sec:eval-baselines}

We compare four baselines that form an ablation ladder. Each rung adds exactly one mechanism on top of the previous.

\textbf{B1 (zero-shot).} The M3A loop on Qwen3-VL-8B-Instruct with no FSM and no $L_C$ in the action-selection prompt. The zero-knowledge floor.

\textbf{B2 (static $L_C$ injection).} The same agent and generation config as B1, but when the target app's Play category is in the source pool the corresponding $L_C$ text is injected between the prefix and the dynamic block of the action-selection prompt. The summary prompt is unchanged. $L_C$ is frozen---no online evolution and no weight update. Tier-C apps receive no injection (the category look-up returns $\bot$) and the prompt degrades byte-identically to B1, giving us a clean null control.

\textbf{B3 (evolved $L_C$).} The FSM population is seeded from the category-matched $L_C$ and grown online by LAYER-2 mutation on $T_{\text{adapt}}$ for 20 iterations ($M{=}2$, $N{=}1$). Claude Opus 4.7 plays the role of the frozen reference model. LoRA is not attached: B3 isolates the contribution of context evolution. The end-of-adaptation champion is frozen and evaluated on $T_{\text{eval}}$.

\textbf{B4 (joint LoRA and $L_C$ adaptation).} The full EvoFSM-RL method of Section~\ref{sec:method}: a shared LoRA adapter is pretrained on the source pool (Phase~1), and both the LoRA adapter and the FSM population are then co-evolved on each target app's $T_{\text{adapt}}$ (Phase~3) starting from this pretrained LoRA, before frozen evaluation on $T_{\text{eval}}$.

All baselines use $\max\_steps\_multiplier = 10$ on top of the per-template default and the same bf16 generation config. The only differences between rows are the contents of $F$ in the action-selection prompt and whether LoRA is updated.

\subsection{Main results}
\label{sec:eval-main}

Table~\ref{tab:main} summarizes the aggregate $T_{\text{eval}}$ success rate per baseline and per tier. B1 and the B1-paired B2 share the seed list $\{30, 31, 32\}$; the B3-paired B2 was rerun on the canonical $T_{\text{eval}}$ seeds $\{40, 41, 42\}$ so that B3 vs.\ B2 is computed on identical seeds. The two B2 numbers (56.5\% and 60.2\%) bound the seed-set noise floor at $n{=}54$.

\begin{table}[t]
\small
\centering
\begin{tabular}{lccc}
\hline
\textbf{Arm} & \textbf{Tier-B} & \textbf{Tier-C} & \textbf{Overall} \\
             & ($n{=}54$) & ($n{=}51$) & ($n{=}105$) \\
\hline
B1 (zero-shot)              & 47.2\%  & 29.4\%  & 38.6\%  \\
B2 paired with B1           & 56.5\%  & 29.4\%  & 43.3\%  \\
B2 paired with B3           & 60.2\%  & 29.4\%  & 45.2\%  \\
B3 (evolved $L_C$)          & \textbf{63.9\%}  & 31.4\%  & \textbf{48.1\%}  \\
B4 (joint LoRA + $L_C$)     & --      & --      & --      \\
\hline
\end{tabular}
\caption{Aggregate $T_{\text{eval}}$ success rate per arm at $K{=}3$ seeds. Tier-C is the built-in null control for B2 and B3, where the prompt is byte-identical to B1.}
\label{tab:main}
\end{table}

Reading off the ablation deltas (Table~\ref{tab:deltas}):

\begin{table}[t]
\small
\centering
\begin{tabular}{lcc}
\hline
\textbf{Comparison} & \textbf{Tier-B $\Delta$} & \textbf{Tier-C $\Delta$} \\
\hline
B2 $-$ B1   & $+9.3$~pp & $+0.0$~pp \\
B3 $-$ B2   & $+3.7$~pp & $+2.0$~pp \\
B4 $-$ B3   & --        & --        \\
\hline
\end{tabular}
\caption{Ablation deltas on $T_{\text{eval}}$ success rate. The Tier-C column is the null-control panel: B2 and B3 inject nothing on Tier-C, so any non-zero delta there reflects paired-run variance rather than a methodological effect.}
\label{tab:deltas}
\end{table}

\textbf{B2 $-$ B1: static category knowledge transfers, with a clean null control.} On Tier-B, $L_C$ injection lifts mean success rate by $+9.3$~pp. The Tier-C delta is exactly $0.0$~pp---the category look-up returns $\bot$ on Tier-C, the prompts are byte-identical, and the two runs converge. This confirms the Tier-B gain is attributable to $L_C$ content rather than to incidental run conditions.

\textbf{B3 $-$ B2: online evolution narrows the residual gap.} A further $+3.7$~pp on Tier-B comes from letting LAYER~2 evolve on the target app's adaptation set. The Tier-C $+2.0$~pp delta is within the paired-run noise floor; under the null hypothesis (Tier-C reuses the no-injection path on both arms) the expected delta is exactly zero.

\subsection{Per-app breakdown and where each mechanism helps}
\label{sec:eval-analysis}

Aggregates can hide structure. Table~\ref{tab:perapp} gives per-app success rates on Tier-B; each app contributes proportional to its T\_eval template count.

\begin{table}[t]
\small
\centering
\begin{tabular}{l@{\,}c@{\,}c@{\,}c@{\,}c@{\,}c}
\hline
\textbf{App}             & \textbf{n} & \textbf{B1} & \textbf{B2$_{\text{B1}}$} & \textbf{B2$_{\text{B3}}$} & \textbf{B3} \\
\hline
simple\_calendar\_pro    & 18 & 38.9 & 38.9 & 50.0 & 50.0 \\
system\_settings         & 18 & 75.0 & 88.9 & 80.6 & \textbf{91.7} \\
pro\_expense             &  9 & 33.3 & \textbf{66.7} & 66.7 & 66.7 \\
retro\_music             &  3 & 66.7 & 50.0 & 33.3 & \textbf{66.7} \\
camera                   &  3 &  0.0 &  0.0 & 66.7 & 33.3 \\
chrome                   &  3 &  0.0 &  0.0 &  0.0 &  0.0 \\
\hline
Tier-B overall           & 54 & 47.2 & 56.5 & 60.2 & \textbf{63.9} \\
\hline
\end{tabular}
\caption{Per-app Tier-B success rate (\%) on $T_{\text{eval}}$. $n$ counts episodes (templates $\times$ K=3). B2$_{\text{B1}}$ uses seeds $\{30, 31, 32\}$; B2$_{\text{B3}}$ and B3 use seeds $\{40, 41, 42\}$.}
\label{tab:perapp}
\end{table}

Three observations come out of this table.

\textbf{B2 helps most when the category $L_C$ has strong cross-source structural overlap with the target.} The biggest B2 gain over B1 is \texttt{pro\_expense} ($+33.3$~pp, Finance $L_C$ from \texttt{bluecoins}): both apps share nearly identical add/edit/delete-transaction workflows, so the abstract category transfers with little adaptation. \texttt{system\_settings} ($+13.9$~pp, Tools $L_C$ from \texttt{calculator}/\texttt{clock}/\texttt{files}) shows the same pattern. Apps whose category structure diverges visually from the source pool (\texttt{retro\_music} vs.\ source music apps, \texttt{simple\_calendar\_pro} vs.\ list-style productivity apps) do not benefit from static injection alone.

\textbf{B3 helps on the apps that have both signal and budget.} \texttt{system\_settings} carries the headline gain ($+11.1$~pp at $n{=}18$); the population there received 14 successful mutations against a uniformly-rated set of strong starting points. Where B3 does not help---\texttt{simple\_calendar\_pro} ties at 50\%, \texttt{chrome} and \texttt{camera} tie at the floor---the failure modes are structural rather than knowledge-shaped, and we discuss them next.

\textbf{Failure modes are interpretable.} \texttt{simple\_calendar\_pro}'s most prominent T\_eval failure is a month-view title-truncation bug that LAYER~2 \emph{could} express as a verification rule, but the mutation operator did not converge on the right rule within the 20-iteration budget. \texttt{chrome/BrowserMultiply} triggers an upstream parsing issue where the agent emits a valid but mislocated \texttt{open\_app: ``File Manager''} action repeatedly until step exhaustion---a benchmark instrumentation gap that no FSM edit can fix. \texttt{camera/CameraTakeVideo} requires precise gesture timing that no LAYER~2 prose can describe. These are honest examples of where the framework's reach ends, not artifacts of poor execution.

\textbf{Tier-C remains the open frontier.} Across B1, B2, and B3, Tier-C success rate hovers in the 29--31\% band---the three methods are methodologically indistinguishable on Tier-C because none of them have any cross-category knowledge to inject. B4 is the first arm whose mechanism actually exercises a learned signal on Tier-C, and whether weight co-adaptation alone can carry far-transfer performance is the most pointed open question this paper sets up.

\subsection{Statistical caveats}
\label{sec:eval-stat}

At $K{=}3$ our $n$ per arm is 54 on Tier-B and 51 on Tier-C. The B2$-$B1 Tier-B two-sample $z$ is $+0.97$ ($p \approx 0.17$ one-sided); B3$-$B2 Tier-B is $z \approx 0.4$. The directions are consistent across two independently seeded comparisons (paired-B1 and paired-B3 against B2) and the signal concentrates on the largest-$n$ Tier-B app in a methodologically interpretable way, but $n{=}54$ is the formal bottleneck. We treat the B1$\to$B2$\to$B3 trajectory as a directional sequence of effect-size estimates rather than as a sequence of significant claims, and we reserve a wider $K_{\text{eval}}{=}5$ sweep (raising the per-arm $n$ to 90) for the camera-ready. The B4 row is the comparison most likely to clear noise, and it is the outstanding piece of evidence this paper sets up.
